$a\leq b$, $c < d$なる非負整数$a,b,c,d$が存在して
\[
n = 10^{a} + 10^{b} = 2^c + 2^d
\]
となるので, 右の方程式を考察する.

\noindent \ybox{\textbf{Case 1. ($a < b$の場合)}}

$v_2(10^{a} + 10^{b}) = a$より, $v_2(2^c + 2^d) = c = a$となる. よって
\[
10^{a} - 2^a = 2^d - 10^{b}
\]
特に, $2^d - 10^b > 0$より$d>b$だから $v_2(2^d - 10^b) = b$. さらに$v_2$に対するLTEの補題より
\bealn{
v_{2}(10^a - 2^a) &= a + v_2(5^a - 1) \\
&= a + (2 + v_2(a))
}
だから$b = a + 2 + v_2(a)$. よって
\[
10^{a}(1 + 10^{2 + v_2(a)}) = 2^a(1 + 2^{d-a})
\]
よって$v_{5}(1 + 2^{d-a}) = a$かつ, $d-a\equiv 2\pmod{4}$かつ, LTEの補題より
\bealn{
a &= v_5(4^{(d-a)/2} - (-1)^{(d-a)/2})\\
  &= 1 + v_5\parena{\frac{d-a}{2}} \\
  &\leq \frac{d-a}{2}
}
より
\[
\frac{d-a}{2} \geq 5^{a}
\]
なので, $d\geq 2\cdot 5^{a} + a$であり,
\[
2^{d} = 10^{a}(1 + 10^{2 + v_2(a)}) - 2^a \geq 2^{(2\cdot 5^{a} + a)}
\]
一方で
\[
10^{a}(1 + 10^{2+v_2(a)}) < (8\sqrt{2})^{a}\cdot (8\sqrt{2})^{2+a} = 2^{7a + 7}
\]
なので$7a + 7 \geq 2\cdot 5^a + a$の必要があり, これを満たすのは$a=1$に限るが, このとき$d$が整数に定まらないため不適.

\noindent\ybox{\textbf{Case 2. ($a=b$の場合)}}
\[2\cdot 10^{a} = 2^{c} + 2^d\]
である. 両辺の$v_2$を見て$a+1 = c$. よって
\[
2^d = 2(10^{a} - 2^a)
\]
となる. $v_2(10^{a} - 2^a) = a+2+v_2(a)$であったから$d = a+3+v_2(a)$でなければならない. よって上を$2^{a+1}$で割って
\[
2^{2 + v_2(a)} = 5^a - 1
\]
となるので, $v_2(a) \leq a-1$より
\[
5^a = 1 + 2^{2 + v_2(a)} \leq 1 + 2^{a+1} \leq 1 + 4\cdot 5^{a-1} \leq 5^{a}
\]
だから$2^{a+1} = 4\cdot 5^{a-1}$. よって$a=1$を得る(mod 5). これにより
\[
(a,b,c,d) = (1,1,2,4)
\]
が得られ,
\[
n = 10^{1} + 10^{1} = 2^{4} + 2^{2} = \boldsymbol{20}
\]
が解である. 
